% Part: normal-modal-logic
% Chapter: completeness
% Section: complete-consistent-sets

\documentclass[../../../include/open-logic-section]{subfiles}

\begin{document}

\olfileid{nml}{com}{ccs}

\olsection{సంపూర్ణ $\Sigma$-అవిరుద్ధ సమితులు}

$\Sigma$ మోడల్ \tetoken{సూత్రాల}{!!{formula}s} సమితి అనుకుందాం---వాటిని
నార్మల్ మోడల్ తర్కపు స్వీకృతాలుగా లేదా నిర్వచక సూత్రాలుగా భావించండి.
సమితి~$\Gamma$, $\Gamma \Proves/[\Sigma] \lfalse$ అయినప్పుడే
$\Sigma$-అవిరుద్ధం; అంటే, ప్రతి $!A_i \in \Gamma$ అయ్యే
$!A_1 \lif (!A_2 \lif \cdots (!A_n \lif
\lfalse)\dots)$కు $\Sigma$ నుంచి
\tetoken{వ్యుత్పత్తి}{!!{derivation}} లేకపోవడం. ప్రతి లోకాన్ని
ఒక ప్రత్యేక $\Sigma$-అవిరుద్ధ సమితిగా తీసుకునే ``కానానికల్''
నమూనాను నిర్మిస్తాం. ఆ సమితి కేవలం $\Sigma$-అవిరుద్ధమేగాక,
ప్రతి మోడల్ \tetoken{సూత్రం}{!!{formula}} సత్యమూల్యాన్ని
నిర్ణయించే అర్థంలో గరిష్ఠ అవిరుద్ధ సమితి: ప్రతి $!A$కు
$!A \in \Gamma$ లేదా $\lnot !A \in \Gamma$:

\begin{defn}
  సమితి~$\Gamma$ $\Sigma$-అవిరుద్ధమై, ప్రతి~$!A$కు
  $!A \in \Gamma$ లేదా $\lnot !A \in \Gamma$ అయితే,
  అప్పుడు మాత్రమే అది \emph{సంపూర్ణ $\Sigma$-అవిరుద్ధ} సమితి.
\end{defn}

సంపూర్ణ $\Sigma$-అవిరుద్ధ సమితులు~$\Gamma$కు అనేక ఉపయోగకర
ధర్మాలు ఉన్నాయి. మొదట, అవి నిగమన పరంగా సంవృతం: $\Gamma
\Proves[\Sigma] !A$ అయితే $!A \in \Gamma$. ముఖ్యంగా,
\tetoken{ఒక సూత్రం}{!!a{formula}}~$!A \in \Sigma$
అయితే, దాని ప్రతి నిదర్శనం కూడా $\in \Gamma$ అవుతుంది. అంతేకాక,
$\Gamma$లో మూలకత్వం ప్రతిజ్ఞావాక్య సంయోజకాల సత్య షరతులను
ప్రతిబింబిస్తుంది. ``కానానికల్ నమూనా''ను నిర్వచించేటప్పుడు
ఇది ముఖ్యమవుతుంది.

\begin{prop}\ollabel{prop:ccs-properties}
  $\Gamma$ సంపూర్ణ $\Sigma$-అవిరుద్ధ సమితి అనుకుందాం. అప్పుడు:
  \begin{enumerate}
  \item \ollabel{prop:ccs-closed}%
    $\Gamma$,~$\Sigma$లో నిగమన పరంగా సంవృతం.
  \item \ollabel{prop:ccs-sigma}%
    $\Sigma \subseteq \Gamma$.
  \tagitem{prvFalse}{\ollabel{prop:ccs-lfalse}%
    $\lfalse \notin \Gamma$}{}
  \tagitem{prvTrue}{\ollabel{prop:ccs-ltrue}%
    $\ltrue \in \Gamma$}{}
  \item \ollabel{prop:ccs-lnot}%
    $\lnot!A \in \Gamma$ అయితే, అప్పుడు మాత్రమే $!A \notin
    \Gamma$.
  \tagitem{prvAnd}{\ollabel{prop:ccs-land}%
    $!A \land !B \in \Gamma$ అయితే, అప్పుడు మాత్రమే $!A \in \Gamma$
    మరియు $!B \in \Gamma$}{}
  \tagitem{prvOr}{\ollabel{prop:ccs-lor}%
    $!A \lor !B \in \Gamma$ అయితే, అప్పుడు మాత్రమే $!A \in \Gamma$
    లేదా $!B \in \Gamma$}{}
  \tagitem{prvIf}{\ollabel{prop:ccs-lif}%
    $!A \lif !B \in \Gamma$ అయితే, అప్పుడు మాత్రమే $!A \notin \Gamma$
    లేదా $!B \in \Gamma$}{}
  \tagitem{prvIff}{\ollabel{prop:ccs-liff}%
    $!A \liff !B \in \Gamma$ అయితే, అప్పుడు మాత్రమే $!A \in \Gamma$
    మరియు $!B \in \Gamma$ రెండూ, లేదా $!A \notin \Gamma$ మరియు
    $!B \notin \Gamma$ రెండూ}{}
  \end{enumerate}
\end{prop}

\begin{proof}
  \begin{enumerate}
  \item $\Gamma \Proves[\Sigma] !A$ అయినా $!A \notin
    \Gamma$ అనుకుందాం. $\Gamma$ సంపూర్ణ $\Sigma$-అవిరుద్ధం
    కనుక $\lnot!A \in \Gamma$. అప్పుడు $\Gamma$ విరుద్ధమవుతుంది;
    ఎందుకంటే $!A, \lnot !A \Proves[\Sigma] \lfalse$.

  \item $!A \in \Sigma$ అయితే $\Gamma \Proves[\Sigma] !A$;
    నిగమన సంవృతత వల్ల, అంటే \olref{prop:ccs-closed} అంశం వల్ల,
    $!A \in \Gamma$.

  \tagitem{prvFalse}{$\lfalse \in \Gamma$ అయితే $\Gamma
    \Proves[\Sigma] \lfalse$; కాబట్టి $\Gamma$
    $\Sigma$-విరుద్ధమవుతుంది.}{}

  \tagitem{prvTrue}{$\Gamma \Proves[\Sigma] \ltrue$ కనుక
    నిగమన సంవృతత, అంటే \olref{prop:ccs-closed} అంశం వల్ల,
    $\ltrue \in \Gamma$.}{}

\item $\lnot!A \in \Gamma$ అయితే, అవైరుధ్యం వల్ల $!A \notin
    \Gamma$; అలాగే $!A \notin \Gamma$ అయితే, $\Gamma$
    సంపూర్ణ $\Sigma$-అవిరుద్ధం కనుక $\lnot!A \in \Gamma$.
    \sourcecorrection{OLTENMLCOMCCS-001}{స్థిర మూలంలోని రెండవ
      దిశలో $!A \notin \Gamma$ నుంచి పొరపాటున $!A \in \Gamma$
      అని ఉంది; సంపూర్ణత ఇచ్చేది $\lnot!A \in \Gamma$.
      సిద్ధాంత వాక్యంలో ఈ రూపమే ఉంది.}

  \tagitem{prvAnd}{\iftag{probAnd}{వ్యాయామం.}{
      $!A \land !B \in \Gamma$ అనుకుందాం. $(!A \land !B) \lif !A$
      సర్వసత్య నిదర్శనం కనుక నిగమన సంవృతత, అంటే
      \olref{prop:ccs-closed} వల్ల $!A \in \Gamma$. అలాగే
      $!B \in \Gamma$. మరోవైపు $!A \in \Gamma$, $!B \in
      \Gamma$ రెండూ అనుకుందాం. $!A \lif (!B \lif (!A \land !B))$
      సర్వసత్య నిదర్శనం కనుక నిగమన సంవృతతతో
      $(!A \land !B) \in \Gamma$.}}{}

  \tagitem{prvOr}{\iftag{probOr}{వ్యాయామం.}{
      $!A \lor !B \in \Gamma$ కాగా $!A \notin \Gamma$,
      $!B \notin \Gamma$ అనుకుందాం. $\Gamma$ సంపూర్ణ
      $\Sigma$-అవిరుద్ధం కనుక $\lnot!A \in \Gamma$ మరియు
      $\lnot !B \in \Gamma$. $\lnot !A \lif (\lnot !B \lif
      \lnot (!A \lor !B))$ సర్వసత్య నిదర్శనం కనుక
      $\lnot(!A \lor !B) \in \Gamma$. అప్పుడు $\Gamma$
      $\Sigma$-విరుద్ధమవుతుంది---ఇది విరోధం. మరో దిశలో, ఏదో
      ఒక ఉపసూత్రం సమితిలో ఉంటే, దాని నుంచి వికల్పం అనుసరించే
      సర్వసత్య నిదర్శనం, నిగమన సంవృతత వల్ల వికల్పమూ సమితిలో ఉంటుంది.
      \sourcecorrection{OLTENMLCOMCCS-002}{స్థిర మూలంలోని
        వ్యాయామం కాని శాఖ వికల్ప-మూలకత్వ తుల్యతలో మొదటి దిశనే
        నిరూపించింది. రెండవ దిశను సర్వసత్య వికల్ప ప్రవేశం,
        నిగమన సంవృతతతో ఇక్కడ స్పష్టం చేశాం.}}}{}

  \tagitem{prvIf}{\iftag{probIf}{వ్యాయామం.}{
      $!A \lif !B \in \Gamma$, $!A \in \Gamma$ అనుకుందాం;
      అప్పుడు $\Gamma \Proves[\Sigma] !B$, కనుక నిగమన సంవృతత
      వల్ల $!B \in \Gamma$. విపరీత దిశకు, $!A \lif !B \notin
      \Gamma$ అయితే $\Gamma$ సంపూర్ణ $\Sigma$-అవిరుద్ధం కనుక
      $\lnot (!A \lif !B) \in \Gamma$. $\lnot(!A \lif !B) \lif !A$
      సర్వసత్య నిదర్శనం కనుక నిగమన సంవృతత వల్ల $!A \in
      \Gamma$. అలాగే $\lnot(!A \lif !B) \lif \lnot !B$
      సర్వసత్య నిదర్శనం కనుక $\lnot !B \in \Gamma$.
      $\Gamma$ $\Sigma$-అవిరుద్ధం కనుక $!B \notin \Gamma$.}}{}

  \tagitem{prvIff}{\iftag{probIff}{వ్యాయామం.}{
      $!A \liff !B \in \Gamma$ అనుకుందాం. $!A \in \Gamma$
      అయితే $!B \in \Gamma$; ఎందుకంటే $(!A \liff !B) \lif
      (!A \lif !B)$ సర్వసత్య నిదర్శనం. అలాగే $!B \in
      \Gamma$ అయితే $!A \in \Gamma$. కనుక $!A \in \Gamma$,
      $!B \in \Gamma$ రెండూ, లేదా $!A \in \Gamma$ గానీ
      $!B \in \Gamma$ గానీ కాని పరిస్థితి.

      విపరీత దిశకు, $!A \liff !B \notin \Gamma$ అనుకుందాం.
      \sourcecorrection{OLTENMLCOMCCS-003}{స్థిర మూలంలోని ఈ
        పరికల్పనలో $!A \lif !B \notin \Gamma$ అని ఉంది;
        వెంటనే వాడే నిషేధ-తుల్యతకు అవసరమైనది
        $!A \liff !B \notin \Gamma$.}
      $\Gamma$ సంపూర్ణ $\Sigma$-అవిరుద్ధం కనుక
      $\lnot (!A \liff !B) \in \Gamma$. $\lnot(!A \liff !B) \lif
      (!A \lif \lnot !B)$ సర్వసత్య నిదర్శనం కనుక $!A \in
      \Gamma$ అయితే $\lnot !B \in \Gamma$; $\Gamma$
      $\Sigma$-అవిరుద్ధం కనుక $!B \notin \Gamma$. అలాగే
      $!B \in \Gamma$ అయితే $!A \notin \Gamma$.
      రెండూ సమితిలో లేని పరిస్థితి కూడా అసాధ్యం: సంపూర్ణత వల్ల
      రెండింటి నిషేధాలూ సమితిలో ఉంటాయి; అప్పుడు ప్రతిజ్ఞావాక్య
      సర్వసత్యం, నిగమన సంవృతత ద్వారా తుల్యత సమితిలోకి వస్తుంది.
      ఇది పరికల్పనకు విరోధం.
      \sourcecorrection{OLTENMLCOMCCS-004}{స్థిర మూలం రెండూ
        సమితిలో ఉండలేవని చూపి, రెండూ సమితిలో లేని పరిస్థితినీ
        నిషేధిస్తుంది; కానీ దానికి వాదన ఇవ్వలేదు. సంపూర్ణతతో
        రెండు నిషేధాలు, వాటి నుంచి తుల్యత వ్యుత్పాద్యతను ఇక్కడ
        చేర్చాం.}
      కనుక $!A \in \Gamma$, $!B \in \Gamma$ రెండూ కావు;
      $!A \notin \Gamma$, $!B \notin \Gamma$ రెండూ కావు.}}{}
  \end{enumerate}
\end{proof}

\begin{probtag}{probAnd,probOr,probIf,probIff}
  \olref[nml][com][ccs]{prop:ccs-properties} నిరూపణను పూర్తిచేయండి.
\end{probtag}

\end{document}
